\section{Integration}

% ── Three orthogonal axes ────────────────────────────────────────────────────
\subsection{Three orthogonal axes}
\textit{(Names on ref sheet; the why and the test-double mapping are not.)}
\begin{itemize}
  \item \textcolor{Bittersweet}{Frequency}: \emph{late + one-time} (wait till end) vs.\ \emph{early + frequent} ($\checkmark$ — surfaces incompatibilities one at a time).
  \item \textcolor{Bittersweet}{Granularity}: \emph{big-bang} (all at once) vs.\ \emph{incremental} ($\checkmark$ — manageable failure isolation).
  \item \textcolor{Bittersweet}{Direction}: \emph{top-down} / \emph{bottom-up} / \emph{sandwich}.
\end{itemize}

\paragraph{Direction $\to$ test-double needed}
\begin{itemize}
  \item \textcolor{Bittersweet}{Top-down}: needs \textbf{stubs} (lower-level not yet ready). Surfaces high-level problems early.
  \item \textcolor{Bittersweet}{Bottom-up}: needs \textbf{drivers} (UI/top-level not ready) to invoke the lower components.
  \item \textcolor{Bittersweet}{Sandwich}: do both, meet in the middle.
\end{itemize}

% ── CI/CD ────────────────────────────────────────────────────────────────────
\paragraph{CI vs.\ CD} \textcolor{Bittersweet}{CI}: build + test on every change. \textcolor{Bittersweet}{CD}: \emph{also} deploys to users automatically. \textcolor{ForestGreen}{CI is a prerequisite for CD}, not a synonym.
